% Generated draft: requires mathematical and source review.
\documentclass[11pt]{article}
\usepackage[margin=25mm]{geometry}
\usepackage{fontspec}
\setmainfont{Latin Modern Roman}
\usepackage{amsmath,amssymb,mathtools,unicode-math}
\setmathfont{Latin Modern Math}
\usepackage{enumitem,needspace,xurl,hyperref}
\setlength{\parindent}{0pt}
\setlength{\parskip}{5pt}
\setlength{\emergencystretch}{3em}
\newcommand{\N}{\mathbb{N}}
\newcommand{\Z}{\mathbb{Z}}
\newcommand{\Q}{\mathbb{Q}}
\newcommand{\R}{\mathbb{R}}
\newcommand{\C}{\mathbb{C}}
\newcommand{\F}{\mathbb{F}}
\newcommand{\A}{\mathbb{A}}
\newcommand{\PP}{\mathbb{P}}
\newcommand{\E}{\mathbb{E}}
\newcommand{\eps}{\varepsilon}
\newcommand{\dd}{\,\mathrm{d}}
\begin{document}
\section*{UnsolvedMath ID: 30001294}
\textbf{Source problem code:} catalogue entry
\subsection{Definitions and mathematical statement}
The supplied question concerns a dispersing billiard map with an absorbing hole. The billiard table, dimension, geometric assumptions, admissible holes, and admissible initial distributions are not specified, so the record does not determine a unique fully quantified problem.

The following measurable formulation makes the observables explicit without supplying those missing hypotheses. Let \((M,\mathcal B)\) be a collision phase space and let \(T\colon M\to M\) denote the billiard map on trajectories for which it is defined. Let \(H\in\mathcal B\) be an absorbing hole and let \(\mu_0\) be a nonzero finite positive measure on \(M\). This measure class is a framework for notation; the intended class of initial distributions remains unspecified. The handling of trajectories on which the billiard map is undefined must also be specified before applying this framework.

Using the convention that a point is removed when its position lies in \(H\), including at time zero, define
\[
\tau_H(x)=\inf\{j\in\Z_{\geq 0}:T^j x\in H\},
\qquad \inf\varnothing=+\infty,
\]
and, for \(n\in\Z_{\geq 0}\),
\[
S_n(H)=\{x:\tau_H(x)\geq n\},
\qquad
m_n(H,\mu_0)=\mu_0(S_n(H)).
\]
Thus \(S_0(H)=M\), and for \(n\geq1\), membership in \(S_n(H)\) means avoiding the hole at times \(0,\ldots,n-1\). The unnormalized distribution after \(n\) surviving iterations is
\[
\nu_n^{H,\mu_0}(A)
=\mu_0\bigl(S_n(H)\cap(T^n)^{-1}(A)\bigr),
\qquad A\in\mathcal B.
\]
Its total mass is \(m_n(H,\mu_0)\). When this mass is positive, the corresponding conditional distribution is
\[
\widehat\nu_n^{H,\mu_0}
=\frac{\nu_n^{H,\mu_0}}{m_n(H,\mu_0)}.
\]
These definitions use one explicit observation-time convention, which the source does not prescribe. In particular, the source does not say whether its symbol \(\mu_n\) means the unnormalized distribution \(\nu_n^{H,\mu_0}\), the conditional distribution \(\widehat\nu_n^{H,\mu_0}\), or another convention.

\textbf{Question as supplied.} For a dispersing billiard map, determine how escape, the surviving mass, and the asymptotic behavior of the evolving distributions depend on variation of the hole \(H\) or of the initial distribution \(\mu_0\). In the notation above, this asks about the dependence of the escape-time distribution, \(m_n(H,\mu_0)\), and the relevant measure sequence on those choices. The admissible variations, requested form of dependence, normalization of \(\mu_n\), and meaning of asymptotic behavior must be specified to obtain a precise theorem-level target. No particular continuity, differentiability, decay law, or convergence assertion is supplied.

The background assessment additionally mentions varying the Hamiltonian. That formulation is not identified with varying the hole: whether changes of the underlying dynamics are also intended remains unresolved.
\subsection{Short English statement}
For a dispersing billiard with an absorbing hole, the question asks how changing the hole or the initial distribution changes escape, remaining mass, and long-term distributions. The supplied record leaves the billiard assumptions, allowed variations, and normalization of \(\mu_n\) unspecified; the displayed definitions give an explicit framework for later clarification.
\subsection{Sources}
\begin{itemize}
\item Local UnsolvedMath registry, permanent ID \(30001294\), source code \texttt{catalogue entry}, title \emph{Escape Dynamics Under Varying Billiard Holes}. The supplied record belongs to the Oberwolfach Reports Open Problems collection. Project registry references are \texttt{data/unsolvedmath/problems.json} and \texttt{catalog/problems.json}; this draft uses the supplied complete record and does not claim to have inspected those files.
\item Supplied report citation: \emph{Dynamische Systeme (2010)}, Oberwolfach Reports, DOI \texttt{10.4171/owr/2009/32}, pp.\ 1787--1842. \url{https://doi.org/10.4171/owr/2009/32}. This is also the supplied literature-source link labeled ``OWR report.'' The record gives proposed year 2009 and names no individual proposer. The report title's year and the other year information are retained as supplied.
\item Supplied literature source: ``Demers--Wright--Young on billiards with holes,'' \url{https://doi.org/10.1007/s00220-009-0941-y}. No more detailed bibliographic information or precise supporting result is supplied.
\item Literature assessment dated exactly \texttt{2026-08-21}: the record labels the problem open and states that the cited results do not settle the general parameter dependence, phrased there in terms of the Hamiltonian or initial measure. This is supplied dated curation, not an independently verified literature assessment. The background expressly requests independent verification of time-sensitive status and cited claims.
\item The supplied statement status is \texttt{corrected\_verified}; its accompanying verification describes only enclosing TeX notation in math delimiters without changing mathematical content. This does not establish mathematical correctness or resolve the formulation issues above.
\end{itemize}
\end{document}
